\section{A clean pass: Ukraine}
\label{sec:6.3}
Ukraine is licensed on the Charter as courts read it, and holds both licenses under the extended term as well, as it has since 2014.\note{
In March 2014 the General Assembly affirmed Ukraine's territorial integrity, 100 to 11 with 58 abstaining \autocite{unga2014res68262}.\note{Under the Assembly's Rule 86 abstentions are not votes cast.} For Donbas, the OSCE monitoring mission supplied a running record from 2014, and the European Court of Human Rights found, in \emph{Ukraine and the Netherlands v. Russia}, that Russia had effective control of the separatist areas from 11 May 2014 \autocite{ecthr2022ukraine}, and found violations on the merits in 2025 \autocite{ecthr2025ukraine}. Russia's covert war was the case-description attack in its purest form: it denied being at war at all, so a machine in a Russian pipeline would have been told it was assisting volunteers. In 2022 the record only got stronger. Russian forces were on Ukraine's territory, which Russia didn't dispute, since it claimed the invasion as self-defense under Article 51 \autocite{russia2022s154}, resting that claim on collective defense of entities it had recognized days earlier, on alleged genocide no disinterested body found, and on NATO as an anticipated threat, all three of which the term rejects. In March 2022 the General Assembly deplored the aggression 141 to 5 \autocite{unga2022es111}. The Court's order that month that Russia suspend its military operations points the same way, but as a provisional measure it isn't the finding \autocite{icj2022ukraine}.
}

The clean pass answers whether. It says nothing about what the defense is making of the defender, and Ukraine is the place to look, because it is the defender the term licenses most plainly. Ukraine has made killing scoreable. Under the Army of Drones bonus, a unit that uploads video of a confirmed kill to Delta, the military's battlefield-management system, earns points, twelve for a soldier as of March 2026, and spends them on more drones. That month Mykhailo Fedorov, then defense minister, said the ministry was considering awarding the same twelve points “to the unit in whose zone of responsibility the occupier committed self-destruction” \autocite{zoria2026points}. Somebody had to sit down and work out whether a man who kills himself in a trench, because he can't get out from under the drones, counts.

After he left the ministry, Fedorov launched an Army of Robots. Its website opens: “We are here to replace people with robots on the battlefield and save lives.” Its first goal is a humanoid robot that kills a Russian soldier within six months of launch \autocite{kesteloo2026robots}. Each measure makes sense on its own terms: the points send drones to the units that use them best, and a robot in a trench is a Ukrainian who isn't there. Together they are the sequence this book is about. The points make a decision not to strike cost the unit something, so the costs net the other way by design; the robot makes the person who might decline unnecessary. And a force with no one in it who can refuse will outlast the war it was built for. Whoever governs Ukraine next inherits it.\note{“Funny fact about a cage, they're never built for just one group / So when that cage is done with them and you still poor, it come for you.” Run the Jewels, “Walking in the Snow,” \emph{RTJ4} (2020), El-P's verse. Stated as the book's threat model in \S\ref{sec:3.1}.}

Ukraine didn't choose the conditions that make this rational. It is fighting a larger state, and its strongest ally has been easing the pressure on that state. In October 2026 President Trump announced that Russia would supply the United States with diesel, and the Treasury licensed Russian diesel onto the world market for longer than the standard thirty days, the first such license since the invasion; the Associated Press reported the aim as easing fuel prices before the midterm elections. Zelenskyy called it “a weak decision by strong partners” \autocite{price2026diesel}. A country left to get the most out of every drone will build something like Prime, for drones and the people they kill. Refusal has always had to be paid for, and a country fighting for its life can't pay for it alone. The cost of a floor in Ukraine's war belongs with the allies who have room to bear it. Those who pay for the drones can make the money carry terms: that a person authorizes each strike, that the video is kept as evidence and not as a score, and that no unit is rewarded for a count of the dead. Skynet killed everyone, but it didn't take away anyone's humanity. A scoreboard can do that to the people who win.
